\documentclass[12pt]{article}

\usepackage[utf8]{inputenc}
\usepackage[T1]{fontenc}
\usepackage[brazil]{babel}

\usepackage{geometry}
\geometry{margin=2.5cm}
\usepackage{setspace}
\onehalfspacing

\usepackage{amsmath, amssymb, amsthm}
\usepackage{mathtools}

\numberwithin{equation}{section}

\theoremstyle{plain}
\newtheorem{theorem}{Teorema}[section]
\newtheorem{proposition}[theorem]{Proposição}
\newtheorem{lemma}[theorem]{Lema}
\newtheorem{corollary}[theorem]{Corolário}

\theoremstyle{definition}
\newtheorem{definition}[theorem]{Definição}
\newtheorem{example}[theorem]{Exemplo}
\newtheorem{exercise}{Exercício}[section]

\theoremstyle{remark}
\newtheorem*{remark}{Observação}

\usepackage{hyperref}
\usepackage{csquotes}

\usepackage[
    backend=biber,
    style=authoryear,
    maxcitenames=2
]{biblatex}

\addbibresource{/mnt/c/Users/nicho/Documents/nicholasvoltani.github.io/bibliography.bib}

\newcommand{\R}{\mathbb{R}}
\newcommand{\pref}{\succeq}
\newcommand{\strictpref}{\succ}
\newcommand{\indiff}{\sim}

% ------------------------------------------------
% Title (fill manually)
% ------------------------------------------------
\title{}
\author{}
\date{}

\begin{document}

\maketitle

\begin{longtable}[]{@{}l@{}}
\toprule\noalign{}
\endhead
\bottomrule\noalign{}
\endlastfoot
date: ``2026-03-15'' \\
tags: \\
- daily \\
draft: ``false'' \\
aliases: \\
\end{longtable}

\begin{quote}
``ESG as an acronym for ``environmental, social, governance'' is a common denominator of the discourse using the term, but a deeper examination reveals that little beyond that understanding is fixed. The word that follows the famous refrain of ``environmental, social, governance'' shapeshifts from ``criteria'' to ``factors,'' ``standards,'' ``strategies,'' ``risks,'' ``issues,'' ``activity,'' or even ``goals.'''' \parencite[, p. 3]{Pollman2022}
\end{quote}

No acima, porém, há uma asserção \textit{do meio financeiro} deste acrônimo. No âmbito popular/leigo, porém, ele pode assumir outros tons mais ``moralistas'': \textgreater{} ``Popular use of the term ESG has even seemed to take on some of these normative views or cultureladen notions that transcend technical ideas of investment screens, impact reporting, or the like. In common parlance, one regularly hears things such as ``startups need ESG,'' buying a certain asset class such as gold is ``not very ESG'' or that companies can ``be'' or ``not be'' ESG.'' \parencite[, p. 4]{Pollman2022}

\hypertarget{o-surgimento-do-termo-esg-onu-relatuxf3rio-who-cares-wins-2004}{%
\section{\texorpdfstring{O surgimento do termo ESG: ONU, relatório \textit{Who Cares Wins} (2004)}{O surgimento do termo ESG: ONU, relatório Who Cares Wins (2004)}}\label{o-surgimento-do-termo-esg-onu-relatuxf3rio-who-cares-wins-2004}}

Desde seu surgimento formal, o ESG sempre teve uma substância \textit{financeira}, sempre se tratou sobre interesses \textit{financeiros} e corporativos. Desde o primeiro parágrafo da primeira seção após a introdução, o relatório da ONU ``\textit{Who Cares Wins}'' \parencite{GlobalCompact2004} já emprega o discurso costumeiro deste meio: falar sobre sustentabilidade no sentido comum da sociedade civil, e logo emendando com os interesses de, e retornos devido a, investidores e agentes econômicos. Tal parágrafo lê: \textgreater{} ``Ultimately, successful investment depends on a vibrant economy, which depends on a healthy civil society, which is ultimately dependent on a sustainable planet. {[}Discurso de sociedade civil até aqui.{]} In the long-term, therefore, investment markets have a \textit{clear self-interest} in contributing to better management of environmental and social impacts in a way that contributes to the sustainable development of global society. {[}Mediação para o ''que me interessa'' dos investidores.{]} A better inclusion of environmental, social and corporate governance (ESG) factors in investment decisions will ultimately contribute to \textit{more stable and predictable markets}, which is \textit{in the interest of all market actors}. {[}Um \textit{pitch} para a adoção de tais medidas pelo meio financeiro.{]}'' \parencite[, p. 3, grifo meu]{GlobalCompact2004}

De fato, os \textit{overall goals} deste relatório são todos pertinentes ao mercado financeiro: ``Stronger and more resilient \textit{financial markets}'', ``Contribution to sustainable \textit{development}''\footnote{O que pode ser interpretado como ``um mundo mais sustentável'', mas que sabe-se que quer, realmente, dizer ``desenvolvimento \textit{capitalista}'', sustentável o que seja.}, ``Awareness and mutual understanding of \textit{involved stakeholders}''\footnote{Evidentemente o \textit{share}holder tem a prioridade. Sem ele, essa conversa nem estaria existindo, afinal!}, ``Improved trust in \textit{financial institutions}'' \parencite[, p. v]{GlobalCompact2004}

A própria proposição deste novo termo --- ESG, ``\textit{environmental, social, governance}'' --- é uma escolha deliberada e explicitamente explicada, como uma busca para desvencilhar-se de termos pré-existentes --- p.~ex. \textit{sustentabilidade} --- que já são carregados de outros significados.\footnote{``Throughout this report we have refrained from using terms such as sustainability, corporate citizenship, etc., in order to avoid misunderstandings deriving from different interpretations of these terms. We have preferred to spell out the environmental, social and governance issues which are the topic of this report.'' \parencite[, p. 1-2]{GlobalCompact2004}}

Em particular, a inclusão do termo ``governança'' é também deliberada, pois \textgreater{} ````Sound corporate governance and risk management systems are crucial pre-requisites to successfully implementing policies and measures to address environmental and social challenges. This is why we have chosen to use the term ``environmental, social and governance issues'' throughout this report, as a way of highlighting the fact that these three areas are closely interlinked.'' \parencite[, p. 2]{GlobalCompact2004}

Tal escolha não foi feita no vácuo, é claro, tendo alguma influência dos escândalos corporativos nos anos 2001-2002, em particular pela empresa Enron \parencite[, p. 14]{Pollman2022}.

\hypertarget{o-propuxf3sito-do-esg-segundo-sua-definiuxe7uxe3o}{%
\subsection{O propósito do ESG, segundo sua definição}\label{o-propuxf3sito-do-esg-segundo-sua-definiuxe7uxe3o}}

O próprio relatório da ONU destaca que a adesão de tais princípios ESG está no interesse das empresas --- o aumento do \textit{shareholder value}, claro ---, pois, ao fazê-lo, elas \textgreater{} ``{[}lidam{]} melhor com riscos relacionados a problemas ESG emergentes, {[}antecipam{]} mudanças regulatórias ou padrões de consumo, e {[}acessam{]} novos mercados ou {[}reduzem{]} custos'' \parencite[, p. 9]{GlobalCompact2004}

Além disso, também lidam com questões de riscos reputacionais.

Ao lidar com tais problemas não individualmente, mas como um todo, como o \textit{framework} do ESG prega, ``empresas bem-sucedidas {[}podem alcançar{]} os melhores resultados em termos de criação de valor'' (ibid.).

Aborda, porém, não só riscos a curto prazo, mas com um conceito de ``materialidade'' a mais longo prazo, ``one that includes longer time horizons (10 years and beyond) and intangible aspects impacting company value.'' \parencite[, p. 15]{GlobalCompact2004}.

\hypertarget{as-acepuxe7uxf5es-possuxedveis-do-termo-esg}{%
\section{As acepções possíveis do termo ESG}\label{as-acepuxe7uxf5es-possuxedveis-do-termo-esg}}

Por sua própria criação, ESG não possui um significado bem delimitado. Embora possua um escopo definido \textit{a priori} --- questões de meio ambiente, sociais e de governança ---, não impõe-se quais problemas dentro destes campos de ação estão fora desta categoria, ou melhor, que ``não são ESG''.

Não há uma definição unívoca fornecida pelo relatório da ONU, além de que deveria se tratar de ``um esforço conjunto de instituições financeiras'' para ``desenvolver diretrizes e recomendações sobre como integrar melhor problemas ambientais, sociais e de governança ao gerenciamento de ativos {[}\textit{asset management}{]}, serviços de corretagem de títulos {[}\textit{securities brokerage services}{]} e funções associadas de pesquisa'' \parencite[, p. vii]{GlobalCompact2004} \#\# 1) Como fator de análises de investimento \textgreater{} ``The term ESG has been, and is, often still used in this vein as a way of referring to a set of issues that should be integrated into investment analysis. As a tool, ESG is often broken into component parts of E, S, and G, and explained by reference to underlying content that would be relevant to investor decision-making.'' \parencite[, p. 21]{Pollman2022}

Embora invoque-se, frequentemente, discursos sobre a melhoria da sociedade civil etc., trata-se, em última instância, de remeter questões/problemáticas na esfera do ESG à esfera financeira. Vide Pollman, tal caráter já se faz presente no relatório original: \textgreater{} ``Although the report {[}\textit{Who Cares Wins}{]} sometimes referred to broader goals such as ``contribut{[}ing{]} to the sustainable development of global society,'' invoking language in the spirit of the UN Global Compact, it heavily emphasized the \textit{``business case'' justification and alignment with long-term value for shareholders}. On the whole, the picture that emerges from the report is that ESG refers to ``information,'' ``issues,'' ``factors,'' or ``criteria'' that should be \textit{integrated into ``normal'' and ``mainstream'' investment analysis}.'' \parencite[, p. 20-1, grifo meu]{Pollman2022}

\hypertarget{esg-enquanto-gerenciamento-de-riscos}{%
\subsection{2) ESG enquanto gerenciamento de riscos}\label{esg-enquanto-gerenciamento-de-riscos}}

Fatores ESG podem também ser vantajosos por seu ``impacto potencial em retornos ajustados ao risco a nível de portfólio {[}\textit{portfolio-level risk-adjusted returns}{]}'' \parencite[, p. 22]{Pollman2022}. De fato, quanto a isso, Pollman observa que \textgreater{} ``Numa era em que más relações públicas ou escândalos corporativos podem ter efeitos devastadores no valor da marca {[}\textit{brand}{]} de uma empresa, o envolvimento de partes interessadas, como consumidores e funcionários, através de''práticas ESG'' pode fornecer informações úteis para gerir relacionamentos-chave e mitigar riscos.'' \parencite[, p. 23]{Pollman2022}

Eu diria que, aqui, há concordância com o item anterior: Enquanto ferramenta no arsenal de análises de investimento, ESG frequentemente aparece na forma de \textit{screenings}, i.e.~triagens de fatores que sejam pertinentes a certa análise financeira (objetos, indivíduos, ativos, etc.) sob filtros predeterminados; tais tarefas são comumente delegadas a equipes de \textit{compliance} de empresas corporativas.

Trata-se, portanto, de uma função corporativa que busca gerenciar riscos, ou melhor, mitigar os custos a eles (potencialmente) associados. Neste ínterim, portanto, o ESG assume a forma de ``a set of practices for proactive risk management, whether at the firm or portfolio level'' \parencite[, p. 24]{Pollman2022}.

\hypertarget{esg-enquanto-responsabilidades-sociais-corporativas}{%
\subsection{3) ESG enquanto ``responsabilidades sociais corporativas''}\label{esg-enquanto-responsabilidades-sociais-corporativas}}

Há interpretações diferentes sobre a relação entre ESG e o anterior CSR (\textit{Corporate Social Responsabilities}). Para alguns, ESG não é mais que um sinônimo; para outros, ESG trata-se de \textbf{mensurar} o cumprimento de tais responsabilidades. Sucintamente colocado, este último ponto pode ser colocado como: \textgreater{} ``CSR is adherence to the actual values of corporate stakeholders, and ESG is a set of measurements from which conclusions about CSR can be drawn.'' (LoPucki \textit{apud} Pollman, 2022, p.~25)

Nesta acepção, ESG assume um caráter \textit{normativo}, de valores éticos, buscando ``injetar consciência social nas decisões corporativas e de investimentos individuais'' (Larcker \textit{et al} \textit{apud} Pollman, 2022, p.~25).

\begin{quote}
``For some, this usage may stem from a nuanced understanding or belief that \textit{broad social benefits may flow from using ESG} as a tool for enhanced investment analysis \parencite[, p. 24, grifo meu]{Pollman2022}''
\end{quote}

\hypertarget{esg-como-preferuxeancia}{%
\subsection{4) ESG como preferência}\label{esg-como-preferuxeancia}}

Tais fatores ESG também podem ser vistos --- não tão diferentemente do caso acima --- como adaptação do meio corporativo às mudanças da sociedade e do mercado.

\begin{quote}
``In this spirit, investors and a wide range of stakeholders seek to align their activities with an expression of their values, whether political, ethical, or social, and ESG is a label vaguely signifying some level of attention to issues beyond the purely financial. \parencite[, p. 27]{Pollman2022}''
\end{quote}

É neste sentido, enfatiza Pollman, que pode-se dizer que uma empresa ``é/não é ESG''\footnote{No mercado, diz-se mais comumente ``ESG \textit{(non-)compliant}'', eu diria.}.

\hypertarget{esg-enquanto-termo-guarda-chuva}{%
\section{ESG enquanto termo guarda-chuva}\label{esg-enquanto-termo-guarda-chuva}}

\hypertarget{esg-como-framework-globalmente-aplicuxe1vel-e-adaptuxe1vel}{%
\subsection{1) ESG como framework globalmente aplicável e adaptável}\label{esg-como-framework-globalmente-aplicuxe1vel-e-adaptuxe1vel}}

É pertinente que este framework seja amplo o suficiente para que abarque o máximo de atividades (do meio financeiro) possível. Ou seja, qualquer diretriz predeterminada pode arriscar de excluir certas áreas.

Por exemplo, a questão de \textit{governança} pode abarcar, em países em desenvolvimento, problemáticas relacionadas a corrupção, e, em países desenvolvidos, problemáticas envolvendo diversidade de gênero etc.

\hypertarget{esg-como-framework-sem-conotauxe7uxf5es-carregadas}{%
\subsection{2) ESG como framework sem conotações carregadas}\label{esg-como-framework-sem-conotauxe7uxf5es-carregadas}}

Não só tendo como alicerce termos amplos, como também podendo ser moldado pelos próprios investidores em suas respectivas áreas de foco.

\hypertarget{esg-como-framework-que-abrange-uma-ampla-gama-de-investidores-no-que-tange-a-seu-alcance-no-mercado}{%
\subsection{3) ESG como framework que abrange uma ampla gama de investidores --- no que tange a seu alcance no mercado}\label{esg-como-framework-que-abrange-uma-ampla-gama-de-investidores-no-que-tange-a-seu-alcance-no-mercado}}

\begin{quote}
``The \textit{Who Cares Wins} initiative explicitly framed ESG in terms of the business case for integrating issues into mainstream investment analysis, chose a term that was facially more neutral than other existing terms, interjected ``governance'' which had widespread buy-in from mainstream market actors, and emphasized the theme of aligning goals between those of the financial industry and the UN. \parencite[, p. 30]{Pollman2022}''
\end{quote}

A questão de ESG tornou-se aderente ao mercado financeiro pois conseguiu integrar ``leis, mercados e culturas orientadas a \textit{shareholders}'' ``na já existente''máquina de governança corporativa''\,'' \parencite[, p. 30]{Pollman2022}.

Permitiu-se, portanto, sua compreensão pelo mercado financeiro como algo ``\textit{value enhancing}'' --- ao contrário da percepção de CSR pelo mercado ---, criando ````oportunidades de negócios'' para uma ampla gama de \textit{players} institucionais como \textit{asset managers}, agências de \textit{ratings}, firmas de contabilidade etc.'' \parencite[, p. 30-1]{Pollman2022}

Neste ínterim, \textgreater{} ``Creating a term that could present itself as neutral or value-enhancing, while at the same time welcoming proponents of previous ``social''-related concepts, enabled a diverse group of investors and stakeholders to embrace activity under such a term.'' \parencite[, p. 31]{Pollman2022}

\hypertarget{arbitragem-de-sustentabilidade}{%
\section{Arbitragem de sustentabilidade}\label{arbitragem-de-sustentabilidade}}

\begin{quote}
``Without an integrated approach to ESG factors, ``sustainability arbitrage'' is possible for both companies and investors. Good performance on one issue, such as low-carbon product development, could be strategically used to mask another, such as poor labor practices.'' \parencite[, p. 34]{Pollman2022}
\end{quote}

Tal fator realmente ocorre, como relatado por \textcite{Simpson2021}, em que uma quantidade substancial das empresas incluídas na rubrica ESG da MSCI tiveram melhoras em seus \textit{ratings} ESG simplesmente por mudanças nas diretrizes ESG. Embora, neste caso, houve mudanças nos \textit{ratings} meramente por mudanças ``exógenas'' às empresas, também pode ocorrer ações ``ESG'' delas que possam ``compensar'' por ações ``não-ESG''.

Neste mesmo ínterim, naturalmente pode ocorrer situações em que o avanço em uma esfera do ESG pode ser detrimental a outra. \textgreater{} ``{[}a{]}dverse employment impacts are to be expected in companies in certain sectors such as energy and some regions that will have to execute an extensive transformation to reduce their GHG emissions and to ultimately stay on a path consistent with the net zero ambitions.'' Environmental concerns and labor interests ``are not always reconcilable'' and divesting or decommissioning brown assets or transforming a business to new technology can lead to workers losing relevant skills, having lower wages, or getting laid off. (Gözlügöl \textit{apud} Pollman, 2022, p.~34)''

É neste contexto que surge o conceito de Transição Justa: lidar com conflitos entre a esfera ambiental e a esfera social de uma maneira sustentável.

\hypertarget{cruxedticas-ao-esg}{%
\section{Críticas ao ESG}\label{cruxedticas-ao-esg}}

\begin{quote}
``The critiques of ESG are wide-ranging, from assertions of confusion, unrealistic expectations, and greenwashing to notions that it is crowding out other solutions or inhibiting accountability.'' \parencite[, p. 35]{Pollman2022}
\end{quote}

\begin{center}\rule{0.5\linewidth}{0.5pt}\end{center}

\hypertarget{referuxeancias}{%
\subsubsection{Referências}\label{referuxeancias}}

POLLMAN, Elizabeth. The origins and consequences of the ESG moniker. \textbf{Institute of Law and Economics Research Paper}, n.~22--23, 2022.

SIMPSON, Cam; RATHI, Akshat; KISHAN, Sajiel. The ESG Mirage: Sustainable Investing Is Mostly About Sustaining Corporations. \textbf{Bloomberg}, 10 dez. 2021.

\printbibliography

\end{document}
